\documentclass[12pt]{article}
\usepackage[a4paper,margin=1in]{geometry}
\usepackage{amsmath,amssymb}
\usepackage{array}
\usepackage{booktabs}
\usepackage{hyperref}

\title{Assignment 3 Solutions}
\author{}
\date{}

\begin{document}
\maketitle

\section*{Question 1}

Let the village locations be
\[
[2,5,12,18,20].
\]

For a station at coordinate \(X\), the asymmetric installation cost is
\[
\mathrm{Cost}(X)=\sum_{x_i<X} 3(X-x_i) + \sum_{x_i>X} (x_i-X).
\]

\subsection*{(a) Fill in the table}

\begin{center}
\begin{tabular}{@{}cccc@{}}
\toprule
Station position \(X\) & Cost from the left & Cost from the right & Total cost \\
\midrule
2  & 0   & 47 & 47 \\
5  & 9   & 35 & 44 \\
12 & 51  & 14 & 65 \\
18 & 105 & 2  & 107 \\
20 & 129 & 0  & 129 \\
\bottomrule
\end{tabular}
\end{center}

Therefore, the optimal location is
\[
X=5
\]
with minimum total cost
\[
44.
\]
So, in short, \(X=5\) is the best choice here.

\subsection*{(b) Mathematical proof of the optimal position}

Assume the village coordinates are sorted:
\[
x_1 \le x_2 \le \cdots \le x_N.
\]

Since the station must be built at one of the village coordinates, compare two consecutive candidate positions \(x_k\) and \(x_{k+1}\). Let
\[
d=x_{k+1}-x_k.
\]

When we move the station from \(x_k\) to \(x_{k+1}\):
\begin{itemize}
    \item Each of the \(k\) villages at or to the left of \(x_k\) becomes \(d\) farther away on the west side, so the total increase is \(3kd\).
    \item Each of the \(N-k\) villages to the right becomes \(d\) closer on the east side, so the total decrease is \((N-k)d\).
\end{itemize}

Hence,
\[
\mathrm{Cost}(x_{k+1})-\mathrm{Cost}(x_k)=3kd-(N-k)d=(4k-N)d.
\]

Now:
\begin{itemize}
    \item If \(4k<N\), then \((4k-N)d<0\), so moving right decreases the cost.
    \item If \(4k>N\), then \((4k-N)d>0\), so moving right increases the cost.
    \item The minimum occurs when \(k\) first reaches \(N/4\).
\end{itemize}

Therefore, for the \(3{:}1\) asymmetric ratio, the optimal station must be the first village whose rank is at least
\[
\left\lceil \frac{N}{4} \right\rceil.
\]

This is basically the weighted median for this asymmetric cost model.

For \(N=5\),
\[
\left\lceil \frac{5}{4} \right\rceil = 2,
\]
so the optimal village is the second village in sorted order, namely
\[
x_2=5.
\]

\subsection*{(c) Function \texttt{calculate\_asymmetric\_cost}}

Implemented in \href{run:assignment3_solution.py}{assignment3\_solution.py}.

\subsection*{(d) Function \texttt{solve\_greedy}}

Implemented in \href{run:assignment3_solution.py}{assignment3\_solution.py}.

For the sample input
\begin{verbatim}
villages = [2, 5, 12, 18, 20]
\end{verbatim}
the function returns
\begin{verbatim}
(5, 44)
\end{verbatim}

\section*{Question 2}

\subsection*{(a) DP state}

Define
\[
DP[i][j]
\]
to be the minimum total cost to serve the first \(i\) villages using exactly \(j\) stations.
You can think of this as the ``best answer so far'' after covering the first \(i\) villages with \(j\) stations.

\subsection*{(b) Definition of \(\mathrm{cost}(p,i)\)}

\(\mathrm{cost}(p,i)\) is the minimum cost of serving the contiguous block of villages
\[
x_p, x_{p+1}, \dots, x_i
\]
with exactly one station placed at the best village location in that block.

Formally,
\[
\mathrm{cost}(p,i)
=
\min_{X \in \{x_p,\dots,x_i\}}
\left(
\sum_{x_t<X} 3(X-x_t) + \sum_{x_t>X} (x_t-X)
\right).
\]

\subsection*{(c) Recurrence and base case}

The recurrence is
\[
DP[i][j] = \min_{1 \le p \le i} \left( DP[p-1][j-1] + \mathrm{cost}(p,i) \right).
\]

Base cases:
\[
DP[0][0]=0,
\]
\[
DP[i][0]=\infty \quad \text{for } i>0,
\]
\[
DP[0][j]=0 \quad \text{for } j\ge 1.
\]

\subsection*{(d) Example: villages \([1,2,5,6]\), \(K=2\)}

First compute the one-station block costs. This part is pretty mechanical:
\[
\mathrm{cost}(1,1)=0,\quad
\mathrm{cost}(1,2)=1,\quad
\mathrm{cost}(1,3)=5,\quad
\mathrm{cost}(1,4)=10,
\]
\[
\mathrm{cost}(2,2)=0,\quad
\mathrm{cost}(2,3)=3,\quad
\mathrm{cost}(2,4)=7,
\]
\[
\mathrm{cost}(3,3)=0,\quad
\mathrm{cost}(3,4)=1,\quad
\mathrm{cost}(4,4)=0.
\]

So the DP table is:

\begin{center}
\begin{tabular}{@{}ccccc@{}}
\toprule
Stations \(j\) / Villages \(i\) & 1 & 2 & 3 & 4 \\
\midrule
1 station  & 0 & 1 & 5 & 10 \\
2 stations & 0 & 0 & 1 & 2 \\
\bottomrule
\end{tabular}
\end{center}

Explanation of the second row:
\begin{itemize}
    \item For \(i=1\), the cost is \(0\) because one station serves village \([1]\), and the second station is unused.
    \item For \(i=2\), the cost is \(0\) by splitting into \([1]\) and \([2]\).
    \item For \(i=3\), the cost is \(1\) by splitting into \([1,2]\) and \([5]\), giving \(1+0=1\).
    \item For \(i=4\), the cost is \(2\) by splitting into \([1,2]\) and \([5,6]\), giving \(1+1=2\).
\end{itemize}

Therefore, the optimal solution is:
\begin{itemize}
    \item Place one station at \(1\) for villages \([1,2]\).
    \item Place one station at \(5\) for villages \([5,6]\).
    \item Minimum total cost \(=2\).
\end{itemize}
So the final answer here is quite neat: split the villages into two small nearby groups.

\section*{Code Note}

The Python file also includes a helper function \texttt{solve\_dp(villages, k)} that reconstructs the optimal village blocks for the multi-station case.

\end{document}
